% !TEX TS-program = xelatex
% !TEX encoding = UTF-8 Unicode
% !Mode:: "TeX:UTF-8"

\documentclass{resume}
\usepackage{zh_CN-Adobefonts_external} % Simplified Chinese Support using external fonts (./fonts/zh_CN-Adobe/)
%\usepackage{zh_CN-Adobefonts_internal} % Simplified Chinese Support using system fonts
\usepackage{linespacing_fix} % disable extra space before next section
\usepackage{cite}
\usepackage{ragged2e}
\usepackage{array}
\usepackage{setspace}
\usepackage{xcolor}
\usepackage{arydshln}
\usepackage{xeCJK}
% \usepackage{ulem}
% \usepackage{enumerate}
\usepackage{hyperref}
\hypersetup{
    colorlinks=true,            %链接颜色
    % linkcolor=blue,             %内部链接
    % filecolor=magenta,          %本地文档
    urlcolor=blue,              %网址链接
    % pdftitle={CV of Yang JS},
    % pdfpagemode=FullScreen,
    }

\begin{document}
% \pagenumbering{gobble} % suppress displaying page number
\pagenumbering{arabic}

% \section{\faAddressCard\,  基本信息}
\begin{center}
  \begin{minipage}[t]{0.75\textwidth}
    \vspace{0pt}
    \hspace{-6pt}\name{杨家树 \Large{西安建筑科技大学}}
    \vspace{12pt}
    % \basicInfo{
    %     % \email{yang\_jiashu@outlook.com} \textperiodcentered\ 
    %     % \phone{+86 18717811089} \textperiodcentered\ }
    %     \leftline{\textcolor[rgb]{0,0.44,0.75}{\faEnvelope}\ yang\_jiashu@outlook.com \ \quad
    %     \textcolor[rgb]{0,0.44,0.75}{\faPhone}\ +86 18717811089 \ }}      \vspace{6pt}
    \begin{spacing}{1.5}
    \hspace{-7pt}\begin{tabular}{p{.2\textwidth} p{.45\textwidth} p{.35\textwidth}}
        \textbf{性别}：男 &  \textbf{籍贯}：河北宁晋 & \textbf{出生日期}：1994年6月 \\
        \textbf{民族}：汉族 & \textbf{邮箱}： yang\_jiashu@outlook.com & \textbf{联系电话}：18717811089 
    \end{tabular}
    \end{spacing}
  \end{minipage}
  \begin{minipage}[t]{0.24\textwidth}
    \vspace{0pt}
    % \centering
    \rightline{\includegraphics[width=0.55\textwidth]{resume/2021_pic11.jpg}}
  \end{minipage}
\end{center}

\vspace{-12pt}
\section{\faUniversity\ 工作经历}
\subsection{$\bullet$\;\;\textbf{2023.07至今} \heiti{，西安建筑科技大学土木工程学院，西安}}
\vspace{12pt}

% \section{\textcolor[rgb]{0,0.44,0.75}{\faGraduationCap\  教育背景}}
\section{\faGraduationCap\  教育背景}
% \datedsubsection{\textbf{同济大学}, 上海}{2016.09 至今}
\subsection{$\bullet$\;\;\textbf{2016.09 -- 2023.02} \heiti{，同济大学土木工程学院，上海}}
\begin{spacing}{1.15}
推荐免试进入土木工程学院建筑工程系直接攻读博士学位，专业为土木工程，导师为国家杰出青年基金获得者陈建兵教授，获研究生学历、工学博士学位。\
\end{spacing}
% \datedsubsection{\textbf{哈尔滨工业大学}，哈尔滨}{2012.09 -- 2016.06}
\subsection{$\bullet$\;\;\textbf{2012.09 -- 2016.06}\heiti{，哈尔滨工业大学土木工程学院，哈尔滨}}
\begin{spacing}{1.15}
就读于土木工程学院土木工程专业力学方向，获本科学历、工学学士学位。
\end{spacing}
% \vspace{-2pt}

% \section{\textcolor[rgb]{0,0.44,0.75}{\faPlane\  访学交流}}
\section{\faPlane\  访学交流}
% \datedsubsection{\textbf{莱布尼兹汉诺威大学}, 德国汉诺威}{2021.07 -- 2022.06 }
\subsection{$\bullet$\;\;\textbf{2021.07 -- 2022.07} \heiti{，莱布尼兹汉诺威大学风险与可靠性研究所，德国汉诺威}}
\begin{spacing}{1.15}
在国家留学基金委 (CSC) 资助下联合培养，为期一年，外方导师为欧洲安全与可靠性协会主席、中国外专局高端外国专家\rm{Michael Beer}教授。\
\end{spacing}

\vspace{-12pt}
% \section{ \textcolor[rgb]{0,0.44,0.75}{\faChartBar\, 博士课题研究内容}}
\section{\faChartBar\, 研究方向}
\begin{spacing}{1.15}
\subsection{$\bullet$\;\;结构参数优化和拓扑优化}
\subsection{$\bullet$\;\;结构随机动力学和动力可靠性分析}
\subsection{$\bullet$\;\;基于动力可靠性的结构参数优化和拓扑优化}
\subsection{$\bullet$\;\;结构抗震可靠性分析}
\end{spacing}
\vspace{6pt}

\vspace{-12pt}
% \section{\textcolor[rgb]{0,0.44,0.75}{\faBook\;\; 学术成果}}
\section{\faBook\;\; 学术成果}
\begin{spacing}{1.1}
% \subsection{\textcolor[rgb]{0,0.44,0.75}{\faFileInvoice\,\; \textbf{SCI期刊论文}}}
\subsection{\faFileInvoice\,\; \textbf{SCI}\heiti{期刊论文}} 
\begin{enumerate}[itemsep=5pt,topsep=2pt]
  \item Wan ZQ, Xie YQ, \underline{\textbf{Yang JS}}*, Lyu MZ, Ghafoori E. Accelerated stochastic model updating via decoupled multi-probability density evolution method with augmented change of probability measure. Mechanical Systems and Signal Processing 246: 113896. (IF=8.9, JCR Q1)

  \item Xiao R, Sushama L, Meguid M, Zayed T, \underline{\textbf{Yang JS}}*. (2026) Time-variant reliability analysis of corroded gas pipelines considering dynamic defect evolution and interaction. Reliability Engineering \& System Safety 267: 111899. (IF=11.0, JCR Q1)
  
  \item \underline{\textbf{Yang JS}}, Lyu MZ*, Chen JB, Xue JY. (2025) Structural design optimization under stochastic excitations considering first-passage probability constraint based on dimension-reduced probability density evolution equation. Reliability Engineering \& System Safety 264: 111378. (IF=11.0, JCR Q1)

   \item Chen G, Liu Y, Ma ZF, \underline{\textbf{Yang JS}}, Beer M, Chian SC*, Wei SJ*. (2025) Assessing extent of building damage following an earthquake: Case study of the 2023 Turkey-Syria doublet. npj Natural Hazards 2: 51.

  \item \underline{\textbf{Yang JS}}, Wan ZQ*, Jensen HA. (2025) An efficient strategy for information reuse in probability density evolution method considering large shift of distributions with multiple random variables. Probabilistic Engineering Mechanics 79: 103728. (IF=3.0, JCR Q1)

  \item Lyu MZ, \underline{\textbf{Yang JS}}*, Chen JB, Li J. (2025) High-efficient non-iterative reliability-based design optimization based on the design space virtually conditionalized reliability evaluation method. Reliability Engineering \& System Safety 254: 110646. (IF=9.4, JCR Q1)

  \item \underline{\textbf{Yang JS}}, Chen JB*, Beer M. (2024) Seismic topology optimization considering first-passage probability by incorporating probability density evolution method and bi-directional evolutionary structural optimization. Engineering Structures 314: 118382. (SCI: IF=5.5, JCR Q1)
  
  \item Weng LL, Acevedo CH, \underline{\textbf{Yang JS}} , Valdebenito MA , Faes MGR , Chen JB*. (2024) An approximate decoupled reliability-based design optimization method for efficient design exploration of linear structures under random loads. Computer Methods in Applied Mechanics and Engineering 432: 117312. (IF=6.9, JCR Q1)

  \item Chen G*, \underline{\textbf{Yang JS}}, Wang RH, Li KQ, Liu Y, Beer M. (2024) Response to discussion of “Seismic damage analysis due to near-fault multipulse ground motion”. Earthquake Engineering \& Structural Dynamics 53(2): 861–866. (SCI: IF=4.5, JCR Q1)

  \item Chen G*, \underline{\textbf{Yang JS}}, Wang RH, Li KQ, Liu Y, Beer M. (2023) Seismic damage analysis subjected to near-fault multi-pulse ground motion. Earthquake Engineering \& Structural Dynamics 52(15): 5099-5116. (SCI: IF=4.5, JCR Q1, \textbf{期刊高被引与高浏览量论文})

  \item Chen G*, \underline{\textbf{Yang JS}}, Liu Y, Kitahara T, Beer M. (2023) An energy-frequency parameter for earthquake ground motion intensity measure. Earthquake Engineering \& Structural Dynamics 52(2): 271–284. (SCI: IF=4.060, JCR Q1)

  \item Weng LL, \underline{\textbf{Yang JS}}, Chen JB*, Beer M. (2023) Structural design optimization under dynamic reliability constraints based on probability density evolution method and quantum-inspired optimization algorithm. Probabilistic Engineering Mechanics 74: 103494 (SCI: IF=2.6, JCR Q1)

  \item Feng CX*, Faes M, Broggi M, Dang C, \underline{\textbf{Yang JS}}, Zheng ZB, Beer M, (2023). Application of interval field method to the stability analysis of slopes in presence of uncertainties. Computers and Geotechnics 153: 105060. (SCI: IF=5.218, JCR Q1)

  \item \underline{\textbf{Yang JS}}, Chen JB*, Beer M, Jensen HA. (2022) An efficient approach for dynamic-reliability-based topology optimization of braced frame structures with probability density evolution method. Advances in Engineering Software 173: 103196. (SCI: IF=4.255, JCR Q1)

  \item \underline{\textbf{Yang JS}}, Chen JB*, Jensen HA. (2022) Structural design optimization under dynamic reliability constraints based on the probability density evolution method and highly-efficient sensitivity analysis. Probabilistic Engineering Mechanics 68: 103205. (SCI: IF=3.350, JCR Q1)
  
  \item \underline{\textbf{Yang JS}}, Jensen HA, Chen JB*. (2022) Structural optimization under dynamic reliability constraints utilizing probability density evolution method and metamodels in augmented input space. Structural and Multidisciplinary Optimization 65(4): 107. (SCI: IF=4.542, JCR Q1)
  
  \item Chen JB*, \underline{\textbf{Yang JS}}, Jensen HA. (2020) Structural optimization considering dynamic reliability constraints via probability density evolution method and change of probability measure. Structural and Multidisciplinary Optimization 62(5): 2499-2516. (SCI: IF=3.377, JCR Q1) 
\end{enumerate}
注: 以上期刊影响因子及分区标注均以发表时为准.
\end{spacing}

% \vspace{-18pt}
% \subsection{\textcolor[rgb]{0,0.44,0.75}{\faFile*\,\; \textbf{EI期刊论文}}}
\subsection{\faFile*\,\; \textbf{EI}\heiti{期刊论文}}
\begin{spacing}{1.3}
\begin{enumerate} [itemsep=5pt,topsep=2pt]
  \item 陈建兵*, 翁丽丽, \underline{\textbf{杨家树}}. (2026) 基于概率密度演化-概率测度变换和量子启发式算法的结构动力可靠性优化设计. 振动工程学报 39(1): 239-248. (EI, 中文核心)
  
  \item \underline{\textbf{杨家树}}, 陈建兵*. (2022) 动力可靠度约束下基于概率测度变换-全局收敛移动渐近线法的结构优化设计. 振动工程学报 35(1): 72-81. (EI, 中文核心，\textbf{中国知网高影响力学术论文})
\end{enumerate}
\end{spacing}

% \vspace{-18pt}
% \begin{spacing}{1.1}
% \subsection{\textcolor[rgb]{0,0.44,0.75}{\faFileSignature\, \textbf{在审论文}}}
% \subsection{\faFileSignature\, \heiti{在审论文}}
% \begin{enumerate} [itemsep=5pt,topsep=2pt]
  % \item Xiao R, Sushama L, Meguid M, Zayed T, \underline{\textbf{Yang JS}}*. Time-variant reliability analysis of corroded gas pipelines considering dynamic defect evolution and interaction. Reliability Engineering \& System Safety. (审稿中)
  % \item Wan ZQ, Xie YQ, \underline{\textbf{Yang JS}}*, Lyu MZ, Ghafoori E. Accelerated stochastic model updating via decoupled multi-probability density evolution method with augmented change of probability measure. Mechanical Systems and Signal Processing. (审稿中)
    % \begin{spacing}{1.15}
    % \item 陈建兵, 翁丽丽, \underline{\textbf{杨家树}}. (2023) 基于概率密度演化-概率测度变换和量子启发式算法的结构动力可靠性优化设计. (EI)
% \end{enumerate}
% \end{spacing}

% \vspace{-18pt}
\begin{spacing}{1.1}
% \subsection{\textcolor[rgb]{0,0.44,0.75}{\faUsers\ \textbf{会议论文}}}
\subsection{\faUsers\ \heiti{国际学术会议报告}}
\begin{enumerate}[itemsep=5pt,topsep=2pt]
  \item \underline{\textbf{Yang JS}}, Lyu MZ, Chen JB. (2024) Dynamic-reliability-based optimization of high-dimensional nonlinear structures under stochastic excitations using dimension-reduced probability density evolution equation. The 10th International Workshop on Reliable Engineering Computing (REC 2024). Beijing, China.

  \item \underline{\textbf{Yang JS}}, Ding YQ. (2024) Seismic-reliability-based topology optimization of frame structures considering stochastic ground motions. The 9th International Symposium on Reliability Engineering and Risk Management (ISRERM 2024). Hefei, China. (\textbf{MS14邀请报告})

  \item \underline{\textbf{Yang JS}}, Wan ZQ. (2024) Dynamic-reliability-based design optimization via probability density evolution method and efficient information reuse strategy. 4th International Conference on Vulnerability and Risk Analysis and Management (ICVRAM2024) \& 8th International Symposium on Uncertainty Modelling and Analysis (ISUMA2024). Shanghai, China. (\textbf{MS06邀请报告})

  \item \underline{\textbf{Yang JS}}, Chen JB. (2024) Reliability-based topology optimization of frame structures under seismic excitations by probability density evolution method. 4th International Conference on Vulnerability and Risk Analysis and Management (ICVRAM2024) \& 8th International Symposium on Uncertainty Modelling and Analysis (ISUMA2024). Shanghai, China.

  \item Weng LL, \underline{\textbf{Yang JS}}, Chen JB. (2023) Dynamic reliability-based design optimization using quantum-inspired algorithm and probability density evolution method assisted by change of probability measure strategy. The 14th International Conference on Applications of Statistics and Probability in Civil Engineering (ICASP14). Dublin, Ireland. (全文)

  \item \underline{\textbf{Yang JS}}, Chen JB, Jensen HA. (2022) A gradient-free approach for dynamic-reliability-based design optimization based on the probability density evolution method. The 13th International Conference on Structural Safety \& Reliability (ICOSSAR2021-2022). Shanghai, China. (全文)

  \item \underline{\textbf{Yang JS}}, Chen JB, Jensen HA. (2021) An efficient approach for dynamic-reliability-based design optimization using the probability density evolution method. The 14th World Congress on Structural and Multidisciplinary Optimization (WCSMO14). Boulder, USA.

  \item \underline{\textbf{Yang JS}}, Chen JB, Jensen HA. (2020) A PDEM-based method for structural optimization under dynamic reliability constraints. The 7th International Symposium on Reliability Engineering \& Risk Management (ISRERM2020). Beijing, China. (全文, \textbf{获学生奖学金})

  \item \underline{\textbf{Yang JS}}, Chen JB, Jensen HA. (2019) An integrated method based on the probability density evolution method for dynamic reliability-based design optimization. The 32th KKHTCNN Symposium on Civil Engineering. Daejeon, South Korea. (全文)
  
  \item \underline{\textbf{Yang JS}}, Chen JB, Jensen HA, Li J. (2019) Structural optimization considering dynamic reliability constraints via probability density evolution method. Proceedings of the 13th World Congress on Structural and Multidisciplinary Optimization (WCSMO13), 128-133. Beijing, China. (全文)
\end{enumerate}
\end{spacing}

\begin{spacing}{1.3}
\subsection{\faUsers\ \heiti{国内学术会议报告}}
\begin{enumerate}[itemsep=5pt,topsep=2pt]
  \item \underline{\textbf{杨家树}}, 律梦泽. (2025) 基于降维概率密度演化方程的结构抗震可靠性优化方法. 第四届土木工程计算与仿真学术会议, 杭州.
  \item \underline{\textbf{杨家树}}, 律梦泽. (2024) 白噪声激励下基于降维概率密度演化方程的结构动力可靠性优化设计. 第十四届全国随机振动理论与应用学术会议暨第十三届全国随机动力学学术会议, 西安.

  \item \underline{\textbf{杨家树}}. (2024) 地震作用下基于动力可靠性的结构拓扑优化方法. 第九届结构设计理论与工程应用学术会议, 南京.

  \item \underline{\textbf{杨家树}}, 陈建兵. (2023) 基于概率密度演化理论的结构拓扑优化设计方法. 第十六届结构工程国际研讨会, 天津.

  \item \underline{\textbf{杨家树}}, 陈建兵. (2023) 基于概率密度演化理论的框架结构拓扑优化. 第十三届全国随机振动理论与应用学术会议暨第十一届全国随机动力学学术会议, 大连.

  \item \underline{\textbf{Yang JS}}, Chen JB, Jensen HA. (2020) An efficient method for sensitivity estimation of failure probability based on the probability density evolution method. The 24th Annual Conference of HKSTAM 2020 \& The 16th Shanghai - Hong Kong Forum on Mechanics and Its Application. Hong Kong. (\textbf{获最佳学生报告奖})
  
  \item \underline{\textbf{杨家树}}, 陈建兵. (2020) 基于概率密度演化-概率测度变换的结构优化设计方法. 第十二届全国随机振动理论与应用学术会议暨第九届全国随机动力学学术会议, 重庆.
\end{enumerate}
\end{spacing}

\begin{spacing}{1.1}
\subsection{\faFileCode\;\, \heiti{技术报告与开源程序}}
\vspace{10pt}
\begin{enumerate}[itemsep=5pt,topsep=2pt]
  \item Chen JB, Lyu MZ, \underline{\textbf{Yang JS}}, Sun TT, Luo Y. (2023) MATLAB and Julia toolboxes for physically-driven DR-PDEE. (Code repository: \url{https://github.com/JCERSM/DR-PDEE-MATLAB} \& \url{https://github.com/JCERSM/DR-PDEE-Julia} )
  
  \item \underline{\textbf{Yang JS}}, Chen JB. (2019) MATLAB toolbox for Probability Density Evolution Method (PDEM). (Code repository: \url{https://github.com/Tree-Yang/Process_Oriented_PDEM} \& \url{https://github.com/Tree-Yang/GF-discrepancy-based-point-selection})

\end{enumerate}
\end{spacing}


\section{\faArchive\;\, 科研项目}
% increase linespacing [parsep=0.5ex]
 \begin{spacing}{1.3}
    \begin{itemize}
    \item 国家自然科学基金青年基金项目（C类），52408212，基于抗震可靠性的高层建筑伸臂桁架智能优化布置方法研究，2025.01-2027.12，主持。
    \item 陕西省博士后科研项目特别资助，2023BSHTBZZ39，考虑时空域随机性的结构抗震可靠性拓扑优化，2024.01-2026.07，主持。
    \item 西安建筑科技大学引进人才科研启动项目，1608724027，随机地震动激励下钢筋混凝土框架结构拓扑优化设计方法，2023.07-2026.07，主持。
    \end{itemize}
 \end{spacing}
% \section{\textcolor[rgb]{0,0.44,0.75}{\faAward\;\, 奖励荣誉}}

\section{\faAward\;\, 奖励荣誉}
% increase linespacing [parsep=0.5ex]

\begin{itemize}
 \begin{spacing}{1.3}
  \item 2025年， Earthquake Engineering \& Structural Dynamics （Wiley）高被引（top cited）与高浏览量（top viewed）论文
  \item 2023年， 中国知网《学术精要数据库》高影响力论文
  \item 2021年， 国家留学基金管理委员会，国家建设高水平大学公派研究生项目资助（12个月）
  \item 2020年， 同济大学，黄金枝奖学金一等奖
  \item 2020年， 第16届沪港力学论坛，最佳学生报告奖
  \item 2020年， 第七届国际可靠性工程与风险管理研讨会（ISRERM7），学生奖学金
  \item 2016年， 同济大学，优秀博士新生奖学金
  \item 2015年， 哈尔滨工业大学，厦门万荣绿建集团奖学金
  \item 2014年， 哈尔滨工业大学，菁华力学专项奖学金
  \item 2013年， 哈尔滨工业大学，三好学生
  \item 2013年， 教育部，国家奖学金
 \end{spacing}
\end{itemize}

% \vspace{-18pt}
% \section{ \textcolor[rgb]{0,0.44,0.75}{\faInfo\;\;\, 其他}}
% \section{\faInfo\;\;\, 其他}
% % increase linespacing [parsep=0.5ex]
% \begin{spacing}{1.15}
% \begin{itemize}
%   \item {\heiti{外语水平}}：英语 CET-6、IELTS 6.5，具有良好的英语阅读、交流和写作能力。
%   \item {\heiti{社会活动经历}}：2020\textasciitilde2021年担任同济大学土木工程学院建筑工程系研究生第十一党支部副书记；多次参与国际学术活动的筹备工作，如2019年JCSS工程可靠性、概率模式规范与风险决策讲习班和第13届国际结构安全性与可靠性学术会议等。
% \end{itemize}
%  \end{spacing}

\section{\faInfo\;\;\, 学术兼职}
% increase linespacing [parsep=0.5ex]
\begin{itemize}
\begin{spacing}{1.3}
  \item 中国工程建设标准化协会城市生命线工程安全专业委员会委员（2023年--\quad）
  \item 国际结构和多学科优化学会（ISSMO）会员
  \item Structural and Multidisciplinary Optimization、Computer Methods in Applied Mechanics and Engineering、Reliability Engineering \& System Safety、Mechanical Systems and Signal Processing、 Engineering Structures、 Probabilistic Engineering Mechanics、Chinese Journal of Aeronautics、ASCE-ASME Journal of Risk and Uncertainty in Engineering Systems, Part A: Civil Engineering、 Journal of Building Engineering、 Results in Engineering等国际期刊审稿人
  \item 《世界地震工程》等中文核心期刊审稿人
\end{spacing}
\end{itemize}
\end{document}
